\chapter{Infinite Series and Convergence}

In Chapter~28, we studied sequences: infinite lists of numbers $a_1, a_2, a_3, \dots$, and asked whether their terms eventually settle on a single value. Now we ask what happens when we try to \emph{add them all together}.

At first sight, adding infinitely many numbers sounds impossible. Addition, as you first learned it, is a finite act: you line up two numbers, or ten, or a million, and you combine them. The process \emph{stops}. So what could it mean to write
\[
1 + \frac{1}{2} + \frac{1}{4} + \frac{1}{8} + \frac{1}{16} + \cdots
\]
and ask for ``the sum''? You can never reach the end. There is no last term to add. Yet many such infinite sums settle onto a definite finite number, while others grow without bound or swing back and forth forever. This chapter develops tools to tell which infinite sums settle down and which do not.

These questions have practical weight. Every time you evaluate $e^{0.3}$ on a calculator, or your computer reports $\sin(1)$, or a physicist models the gravitational pull of a ring of mass, an infinite sum is doing the work underneath. Transcendental functions such as $e^x$, $\sin x$, and $\ln(1+x)$ can be written as infinite sums, as the next chapter shows. But none of it makes sense until we are clear about what it means for an infinite sum to have a value at all, so that is where we begin.

\section{What It Means to Add Infinitely Many Things}
\label{ch29:partial-sums}

The most important idea in the chapter is that we never perform an infinite addition. We \emph{can't}. Instead we perform a sequence of ordinary finite additions and watch where they are heading.

\begin{intuitionbox}
Imagine walking toward a wall. First you cover half the remaining distance, then half of what's left, then half again. After each step you have walked a finite, knowable distance. You will never literally arrive, but it is intuitively clear \emph{where} you are arriving: at the wall. The infinite sum is not the act of taking infinitely many steps. It is the \emph{destination} those steps approach.
\end{intuitionbox}

Figure~\ref{ch29:fig-wall} draws this walk on a number line, with a wall placed $1$ unit away.

\begin{figure}[htbp]
\centering
\begin{tikzpicture}[x=11cm,y=1cm]
\fill[curveblue!55] (0,0) rectangle (0.5,0.5);
\fill[curveblue!35] (0.5,0) rectangle (0.75,0.5);
\fill[curveblue!55] (0.75,0) rectangle (0.875,0.5);
\fill[curveblue!35] (0.875,0) rectangle (0.9375,0.5);
\fill[curveblue!55] (0.9375,0) rectangle (0.96875,0.5);
\fill[curveblue!35] (0.96875,0) rectangle (0.984375,0.5);
\draw[thick] (0,0) -- (1,0);
\draw[tangentred, line width=2pt] (1,-0.3) -- (1,1.1);
\node[tangentred, font=\small, anchor=west] at (1.01,0.9) {wall};
\foreach \x/\l in {0/0, 0.5/\tfrac12, 0.75/\tfrac34, 0.875/\tfrac78, 1/1} {
  \draw (\x,0) -- (\x,-0.12);
  \node[font=\small, below] at (\x,-0.12) {$\l$};
}
\node[font=\small] at (0.25,0.25) {$\tfrac12$};
\node[font=\small] at (0.625,0.25) {$\tfrac14$};
\node[font=\small] at (0.8125,0.25) {$\tfrac18$};
\draw[decorate, decoration={brace, amplitude=5pt}] (0,0.65) -- (0.875,0.65)
  node[midway, above=5pt, font=\small] {after $3$ steps: $S_3 = \tfrac12+\tfrac14+\tfrac18 = \tfrac78$};
\end{tikzpicture}
\caption{Walking to a wall $1$ unit away by halving the remaining gap. Each shaded block is one step. After $n$ steps the distance walked is $1 - \tfrac{1}{2^n}$, and the gap left to the wall is $\tfrac{1}{2^n}$, which shrinks toward $0$.}
\label{ch29:fig-wall}
\end{figure}

Now we make this precise. Given an infinite list of numbers $a_1, a_2, a_3, \dots$, which we call the \emph{terms}, we build a new list out of them, the list of \textbf{partial sums} (running totals):
\[
S_1 = a_1, \qquad
S_2 = a_1 + a_2, \qquad
S_3 = a_1 + a_2 + a_3, \qquad \dots, \qquad
S_n = \sum_{k=1}^{n} a_k .
\]
Here $S_n$ is the sum of the first $n$ terms, and $k$ is a counter that runs from $1$ to $n$. Each $S_n$ is a finite sum: ordinary arithmetic. The infinite sum $a_1 + a_2 + a_3 + \cdots$ is called an \textbf{infinite series} (or just a \textbf{series}), and its value is \emph{defined} to be the limit of the sequence of partial sums, if that limit exists.

\begin{definition}[Convergence of a series]
\label{ch29:def-convergence}
Let $a_1, a_2, a_3, \dots$ be a sequence of real numbers and let $S_n = \sum_{k=1}^n a_k$ be its $n$th partial sum. We say the series $\sum_{k=1}^{\infty} a_k$ \textbf{converges} to the number $S$ if
\[
\lim_{n \to \infty} S_n = S,
\]
and we then write $\sum_{k=1}^{\infty} a_k = S$. If the sequence $(S_n)$ has no finite limit, we say the series \textbf{diverges}.
\end{definition}

\begin{happeningbox}
The symbol $\sum_{k=1}^{\infty} a_k$ does not instruct you to perform infinitely many additions. It is shorthand for a \emph{limit of a sequence}: the sequence of finite partial sums. When we write $\sum_{k=1}^\infty \tfrac{1}{2^k} = 1$, we mean that the sequence of finite sums $\tfrac12, \tfrac34, \tfrac78, \tfrac{15}{16}, \dots$ approaches $1$. An infinite sum is a statement about limits.
\end{happeningbox}

\begin{historybox}
Why was this definition needed? Because without it, even the best mathematicians fell into contradictions.

In the eighteenth century, mathematicians had no precise definition of what an infinite sum meant. They treated infinite series like ordinary finite polynomials, rearranging and grouping terms at will. In 1703, the Italian monk and mathematician Guido Grandi analyzed the alternating series:
\[
  1 - 1 + 1 - 1 + 1 - 1 + \cdots
\]
Grouping adjacent pairs as $(1 - 1) + (1 - 1) + \cdots$ gave $0 + 0 + \cdots = 0$, while shifting the grouping by one term as $1 - (1 - 1) - (1 - 1) - \cdots$ gave $1 - 0 - 0 - \cdots = 1$. Grandi argued that since the sum was poised between $0$ and $1$, its true value must be the average, $1/2$, and claimed this proved God created the universe out of nothing ($0 = \tfrac12 - \tfrac12$). Even Leonhard Euler accepted that $1 - 1 + 1 - \cdots = 1/2$, by putting $x = 1$ into $\frac{1}{1+x} = 1 - x + x^2 - \cdots$. But this kind of purely symbolic algebra, applied without asking whether the sums had values, led to absurd results, such as $1 + 2 + 4 + 8 + \cdots = -1$ from putting $x = 2$ into $\frac{1}{1-x} = 1 + x + x^2 + \cdots$. Critics attacked the new calculus as logically suspect.

Order was restored in 1821 when Augustin-Louis Cauchy published his \emph{Cours d'Analyse}. Cauchy established Definition~\ref{ch29:def-convergence}: an infinite sum has a value if and only if the sequence of its finite partial sums $S_n$ approaches a single, finite limit. Because the partial sums of Grandi's series bounce endlessly between $1$ and $0$ ($1, 0, 1, 0, \dots$), no limit exists, and the series diverges. Cauchy's definition replaced arguments about what such sums ``ought'' to equal with a test that can be checked.

Cauchy's framework made it possible to put mathematical physics and Fourier series (sums of sines and cosines used to describe waves) on firm ground. Yet infinite series still pose open questions. Euler discovered that $\sum_{n=1}^\infty \frac{1}{n^2} = \frac{\pi^2}{6}$, but no closed-form formula (one in terms of familiar constants such as $\pi$) is known for sums with odd exponents, such as $\sum_{n=1}^\infty \frac{1}{n^3}$ or $\sum_{n=1}^\infty \frac{1}{n^5}$. Roger Ap\'ery proved in 1978 that $\sum 1/n^3$ is irrational, but that is all. Similarly, the gap between $1 + \frac12 + \cdots + \frac1n$ and $\ln n$ approaches the Euler--Mascheroni constant $\gamma = \lim_{n\to\infty} \big(\sum_{k=1}^n \frac{1}{k} - \ln n\big) \approx 0.5772$ from Chapter~28, and nobody knows whether $\gamma$ is rational or irrational.
\end{historybox}

So ``convergence'' and ``divergence'' are not new ideas. They are the limit of a sequence from Chapter~28, applied to the particular sequence $(S_n)$ of partial sums. Every theorem about series is a statement about that sequence, and keeping this in mind explains why each convergence test works. Figure~\ref{ch29:fig-partialsums} plots the partial sums of $1 + \tfrac12 + \tfrac14 + \cdots$ (the walk of Figure~\ref{ch29:fig-wall} with every step doubled, so the wall is at $2$).

\begin{figure}[H]
\centering
\begin{tikzpicture}
\begin{axis}[
    width=12cm, height=6cm,
    xlabel={$n$ (number of terms added)}, ylabel={partial sum $S_n$},
    xmin=0, xmax=10.5, ymin=0, ymax=2.2,
    xtick={1,2,3,4,5,6,7,8,9,10},
    samples=10, domain=1:10,
    legend pos=south east, legend cell align=left,
    grid=both, grid style={gray!18},
]
\addplot[only marks, mark=*, mark size=1.6pt, curveblue]
    coordinates {(1,1)(2,1.5)(3,1.75)(4,1.875)(5,1.9375)(6,1.96875)(7,1.984375)(8,1.9921875)(9,1.99609375)(10,1.998046875)};
\addlegendentry{$S_n = \sum_{k=0}^{n-1} 2^{-k}$}
\addplot[tangentred, dashed, thick, domain=0:10.5] {2};
\addlegendentry{limit $= 2$}
\end{axis}
\end{tikzpicture}
\caption{The partial sums of $1 + \tfrac12 + \tfrac14 + \cdots$ climb toward $2$, halving the gap each step but never crossing it. The value of the series is this limit, $2$.}
\label{ch29:fig-partialsums}
\end{figure}

\begin{walkbox}
Before reading on, settle a question in your head on a walk: the terms $\tfrac{1}{2^k}$ get smaller and smaller and shrink toward zero. Is ``the terms shrink to zero'' enough, all by itself, to guarantee the sum is finite? Hold your answer. By the end of Section~\ref{ch29:harmonic} it may surprise you.
\end{walkbox}

\section{The Geometric Series: The One You Can Sum Exactly}
\label{ch29:geometric}

Most convergent series have no simple closed-form sum -- we can often prove a series converges without being able to write its exact total in elementary terms. A major exception that can be evaluated exactly is the \textbf{geometric series}, in which each term is a constant multiple $r$ of the preceding term.

\[
a + ar + ar^2 + ar^3 + \cdots = \sum_{k=0}^{\infty} a\, r^k .
\]

Here $a$ is the first term and the number $r$ is the \emph{ratio}. The sum of $1 + \tfrac12 + \tfrac14 + \cdots$ in Figure~\ref{ch29:fig-partialsums} is the case $a=1$, $r=\tfrac12$. We find the total by computing the partial sums and taking their limit.

\begin{focusbox}
\textbf{The summing trick.} Let $S_n = a + ar + ar^2 + \cdots + ar^{n-1}$ (that's $n$ terms). Multiply the whole thing by $r$:
\[
r S_n = ar + ar^2 + \cdots + ar^{n-1} + ar^{n}.
\]
Now subtract the second line from the first. Almost everything cancels, because both lines contain the terms $ar, ar^2, \dots, ar^{n-1}$; only the $a$ at the start of $S_n$ and the $ar^n$ at the end of $rS_n$ survive.
\[
S_n - rS_n = a - ar^{n} \quad\Longrightarrow\quad (1-r)S_n = a\,(1 - r^{n}).
\]
Provided $r \neq 1$, divide:
\[
\boxed{\,S_n = a\,\frac{1 - r^{n}}{1 - r}\,}.
\]
\end{focusbox}

This is an \emph{exact} formula for the finite partial sum, useful on its own. Now we take the limit. Everything depends on what $r^n$ does as $n \to \infty$. If $\abs{r} < 1$, then $r^n \to 0$ (repeatedly multiplying by something smaller than $1$ in size drives you to zero). If $\abs{r} > 1$, then $\abs{r^n} \to \infty$ and the partial sums blow up. The boundary cases $r = 1$ and $r = -1$ are handled separately in the proof below.

\begin{theorem}[The geometric series]
\label{ch29:thm-geometric}
The geometric series $\sum_{k=0}^{\infty} a\, r^k$ converges if and only if $\abs{r} < 1$, and in that case
\[
\sum_{k=0}^{\infty} a\, r^{k} = \frac{a}{1 - r}.
\]
If $\abs{r} \geq 1$ (and $a \neq 0$) the series diverges.
\end{theorem}

\begin{proof}
From the boxed formula, $S_n = a\dfrac{1 - r^n}{1-r}$ for $r \neq 1$. If $\abs r < 1$ then $r^n \to 0$, so $S_n \to \dfrac{a(1-0)}{1-r} = \dfrac{a}{1-r}$. If $\abs r > 1$ then $\abs{r^n}\to\infty$, so $S_n$ has no finite limit. If $r = 1$, the series is $a + a + a + \cdots$ and $S_n = na \to \pm\infty$. If $r = -1$, the partial sums bounce $a, 0, a, 0, \dots$ and never settle. In every case with $\abs r \ge 1$ the series diverges.
\end{proof}

\begin{whybox}
Why does the condition $\abs{r}<1$ decide everything? Look at the \emph{tail}: the part of the series left over after the first $n$ terms. For a geometric series the tail is itself a geometric series, $ar^n + ar^{n+1}+\cdots = r^n \cdot \frac{a}{1-r}$ (when $\abs r<1$), a copy of the whole series scaled by $r^n$. When $\abs r < 1$ that scaling factor shrinks the leftover tail to nothing, so the partial sums close in on the total. When $\abs r \ge 1$ the terms $ar^n$ do not shrink at all, so the partial sums never settle. Because every tail is a scaled copy of the whole series, we can compute the sum exactly, and only the size of $\abs r$ compared with $1$ matters.
\end{whybox}

\begin{example}[The walk to the wall, finished]
With $a=1$, $r=\tfrac12$: the sum is $\dfrac{1}{1 - \tfrac12} = 2$. Our list started at $k=0$, so this is $1 + \tfrac12 + \tfrac14 + \cdots = 2$, matching Figure~\ref{ch29:fig-partialsums}.
\end{example}

\begin{example}[A repeating decimal is a geometric series]
What is $0.7777\ldots$ really? It is
\[
\frac{7}{10} + \frac{7}{100} + \frac{7}{1000} + \cdots = \sum_{k=1}^{\infty} \frac{7}{10}\left(\frac{1}{10}\right)^{k-1},
\]
a geometric series with first term $a = \tfrac{7}{10}$ and ratio $r = \tfrac{1}{10}$. So its value is
\[
\frac{a}{1-r} = \frac{7/10}{1 - 1/10} = \frac{7/10}{9/10} = \frac{7}{9}.
\]
Every repeating decimal can be written as a geometric series, showing why repeating decimals represent rational numbers. For instance, $0.999\ldots = \sum_{k=1}^\infty \frac{9}{10^k} = \frac{9/10}{1-1/10}=1$.
\end{example}

\begin{connectionbox}
The geometric series is ancient. Around 250~BCE, Archimedes summed $1 + \tfrac14 + \tfrac{1}{16} + \cdots = \tfrac43$ to find the area of a parabolic segment (the region between a parabola and a straight line cutting across it). This was an infinite sum, handled rigorously, two thousand years before ``limit'' had a formal definition. He did not say ``let $n \to \infty$''; instead he proved by his \emph{method of exhaustion} that the area (measured in units of a certain inscribed triangle) could be neither more nor less than $\tfrac43$, squeezing it from both sides. The nineteenth-century definitions of limits and partial sums turned such one-off arguments into a general method that works for any series.
\end{connectionbox}

\begin{recoverybox}
\textbf{Mental check.} What is $\sum_{k=1}^\infty (1/3)^k$? The index begins at $k=1$, so the first term is $a = 1/3$ and the common ratio is $r = 1/3$. By Theorem~\ref{ch29:thm-geometric}, the sum is $\frac{1/3}{1 - 1/3} = \frac{1/3}{2/3} = \frac{1}{2}$.
\end{recoverybox}

\section{Telescoping Series: When the Middle Cancels}
\label{ch29:telescoping}

Besides the geometric series, there is one more family we can sum exactly, and it relies on the same essential move: arrange for the partial sum $S_n$ to collapse to just a few surviving terms. These are the \textbf{telescoping series}, named for the way an old collapsible spyglass folds down to almost nothing.

The idea: suppose each term can be written as a \emph{difference} of consecutive values of some sequence, $a_k = b_k - b_{k+1}$. Then
\[
S_n = \sum_{k=1}^n (b_k - b_{k+1}) = (b_1 - b_2) + (b_2 - b_3) + (b_3 - b_4) + \cdots + (b_n - b_{n+1}).
\]
Every interior term appears twice with opposite signs and cancels. Only the very first and very last survive:
\[
S_n = b_1 - b_{n+1}.
\]
The infinite sum is then $\lim_{n\to\infty} S_n = b_1 - \lim_{n\to\infty} b_{n+1}$, easy as soon as you know where $b_k$ heads. Figure~\ref{ch29:fig-telescope} shows the cancellation.

\begin{figure}[htbp]
\centering
\begin{tikzpicture}[font=\small]
\node (t1) at (0,0)   {$(b_1$};
\node (t2) at (1.0,0) {$-\,b_2)$};
\node at (1.75,0) {$+$};
\node (t3) at (2.5,0) {$(b_2$};
\node (t4) at (3.5,0) {$-\,b_3)$};
\node at (4.25,0) {$+$};
\node (t5) at (5.0,0) {$(b_3$};
\node (t6) at (6.0,0) {$-\,b_4)$};
\node at (6.9,0) {$+$};
\node (dd) at (7.7,0) {$\cdots$};
\node at (8.5,0) {$+$};
\node (t7) at (9.3,0) {$(b_n$};
\node (t8) at (10.5,0) {$-\,b_{n+1})$};
\draw[trapred, thick, bend left=40, <->] (t2.north) to node[above, font=\footnotesize] {cancel} (t3.north);
\draw[trapred, thick, bend left=40, <->] (t4.north) to node[above, font=\footnotesize] {cancel} (t5.north);
\draw[trapred, thick, bend left=40, <->] (t6.north) to node[above, font=\footnotesize] {cancel} (dd.north);
\draw[trapred, thick, bend left=40, <->] (dd.north) to node[above, font=\footnotesize] {cancel} (t7.north);
\draw[targetgreen, thick] (t1.south west) -- (t1.south east);
\draw[targetgreen, thick] (t8.south west) -- (t8.south east);
\node[targetgreen!70!black, below=3pt] at (t1.south) {survives};
\node[targetgreen!70!black, below=3pt] at (t8.south) {survives};
\end{tikzpicture}
\caption{In a telescoping sum each $-b_{k+1}$ is cancelled by the $+b_{k+1}$ that starts the next bracket. Only the first piece $b_1$ and the last piece $-b_{n+1}$ survive, so $S_n = b_1 - b_{n+1}$.}
\label{ch29:fig-telescope}
\end{figure}

\begin{example}[A classic telescoping sum]
Consider $\displaystyle\sum_{k=1}^{\infty} \frac{1}{k(k+1)}$. The terms shrink fast, but is there a sum? Partial fractions (Chapter~24) split each term into a difference:
\[
\frac{1}{k(k+1)} = \frac{1}{k} - \frac{1}{k+1}.
\]
So with $b_k = \tfrac1k$ we have $a_k = b_k - b_{k+1}$, and
\[
S_n = \left(1 - \tfrac12\right) + \left(\tfrac12 - \tfrac13\right) + \cdots + \left(\tfrac1n - \tfrac{1}{n+1}\right) = 1 - \frac{1}{n+1}.
\]
As $n\to\infty$, $\tfrac{1}{n+1}\to 0$, so the series sums to exactly $1$.
\end{example}

\begin{trapbox}
A common error in telescoping series is assuming that every term cancels except the first. To avoid this, write out the partial sum $S_n$ in full and identify which terms remain at \emph{both} ends. Some telescopes leave two survivors on each side (when the difference skips a step, like $b_k - b_{k+2}$), giving $S_n = b_1 + b_2 - b_{n+1} - b_{n+2}$. And the sum exists only if the trailing survivors have a limit. If $b_k$ does not converge, the telescope diverges no matter how much cancels in the middle.
\end{trapbox}

\begin{recoverybox}
\textbf{Mental check.} If $a_k = \frac{1}{k} - \frac{1}{k+2}$, which terms survive in $S_n$? Expand:
\[
S_n = \left(1 - \tfrac13\right) + \left(\tfrac12 - \tfrac14\right) + \left(\tfrac13 - \tfrac15\right) + \left(\tfrac14 - \tfrac16\right) + \cdots
\]
The terms $1$ and $\frac12$ have no subtracted partners, and the last two subtracted terms have no positive partners, leaving
\[
S_n = 1 + \tfrac12 - \tfrac{1}{n+1} - \tfrac{1}{n+2} \;\longrightarrow\; \tfrac32 .
\]
\end{recoverybox}

\section{The First Thing to Check: The $n$th-Term Test}
\label{ch29:nth-term}

Before trying any of the longer tests, there is a quick check that can show at once that a series diverges. It comes from a simple observation about \emph{convergent} series and their terms.

\begin{intuitionbox}
If the partial sums $S_n$ are settling down to a fixed number $S$, then consecutive partial sums must be getting closer and closer to each other. But the gap between consecutive partial sums is exactly one term: $S_n - S_{n-1} = a_n$. So the terms must be shrinking to zero. A convergent sum \emph{cannot} keep adding non-negligible chunks forever -- if it did, it would never settle.
\end{intuitionbox}

\begin{theorem}[The $n$th-term test for divergence]
\label{ch29:thm-nth-term}
If $\sum_{k=1}^{\infty} a_k$ converges, then $\lim_{n\to\infty} a_n = 0$. Equivalently: if $a_n$ does \emph{not} tend to $0$, the series diverges.
\end{theorem}

The second sentence is the \emph{contrapositive} of the first. For any statement ``if $P$ then $Q$'', the contrapositive is ``if $Q$ is false then $P$ is false'', and the two always say the same thing. (``If it is raining, the ground is wet'' has contrapositive ``if the ground is dry, it is not raining''.)

\begin{proof}
Suppose the series converges to $S$, so $S_n \to S$ and also $S_{n-1} \to S$ (same sequence, shifted). Then
\[
a_n = S_n - S_{n-1} \;\longrightarrow\; S - S = 0. \qedhere
\]
\end{proof}

\begin{trapbox}
This theorem is often misused, so read carefully. It says one thing only: \emph{if the terms don't go to zero, the series diverges.} It does \textbf{not} say the converse (the statement with ``if'' and ``then'' swapped). ``The terms go to zero'' does \textbf{not} imply the series converges. That converse is false, and the harmonic series in Section~\ref{ch29:harmonic} is the famous counterexample. The $n$th-term test is a \emph{divergence detector only}. When the terms shrink to zero, the test tells you nothing, and you need another test.
\end{trapbox}

\begin{example}
$\sum \dfrac{n}{2n+1}$ diverges, because $\dfrac{n}{2n+1}\to \tfrac12 \neq 0$. We are adding chunks that stay near $\tfrac12$ forever; no settling is possible. Likewise $\sum (-1)^n$ diverges since the terms don't approach $0$ (they bounce between $1$ and $-1$).
\end{example}

\section{The Harmonic Series: A Genuine Surprise}
\label{ch29:harmonic}

The walking question asked whether terms shrinking to zero is enough for a finite sum. This series answers it:
\[
\sum_{n=1}^{\infty} \frac{1}{n} = 1 + \frac{1}{2} + \frac{1}{3} + \frac{1}{4} + \frac{1}{5} + \cdots
\]
Why is it called the \emph{harmonic} series? The name comes from music: a vibrating string of length $L$ also vibrates in halves, thirds, quarters, and so on, and these overtones (or harmonics) have wavelengths $L/2, L/3, L/4, \dots$. The terms shrink to zero, so the $n$th-term test says nothing against it. Yet the harmonic series \textbf{diverges}: its partial sums grow past every number, though slowly. The proof is short.

\begin{focusbox}
\textbf{Oresme's grouping (c.~1350).} Group the terms in blocks whose lengths double, and bound each block from below:
\[
1 + \frac12 + \underbrace{\frac13 + \frac14}_{> \frac14 + \frac14 = \frac12} + \underbrace{\frac15 + \frac16 + \frac17 + \frac18}_{> \frac18\cdot 4 = \frac12} + \underbrace{\frac19 + \cdots + \frac1{16}}_{> \frac{1}{16}\cdot 8 = \frac12} + \cdots
\]
In each block, replace every term by the \emph{smallest} term in that block; the block is still at least $\tfrac12$. So the sum contains infinitely many chunks, each at least $\tfrac12$: after the first $2^m$ terms, $S_{2^m} \ge 1 + \tfrac{m}{2}$. Add enough blocks and you pass any number you like. The partial sums are unbounded; the series diverges. Figure~\ref{ch29:fig-oresme} draws the blocks.
\end{focusbox}

\begin{figure}[htbp]
\centering
\begin{tikzpicture}
\begin{axis}[
    width=13cm, height=6.5cm,
    xlabel={$n$}, ylabel={term $1/n$},
    xmin=0.3, xmax=16.7, ymin=0, ymax=1.08,
    xtick={1,2,3,4,5,6,7,8,9,10,11,12,13,14,15,16},
    ytick={0.25,0.5,1},
    yticklabels={$\tfrac14$,$\tfrac12$,$1$},
    ticklabel style={font=\small},
    ybar, bar width=0.62, bar shift=0pt,
    axis lines=left, clip=false,
]
\addplot[fill=gray!35, draw=gray!70!black] coordinates {(1,1)(2,0.5)};
\addplot[fill=curveblue!45, draw=curveblue] coordinates {(3,0.3333)(4,0.25)};
\addplot[fill=quotientorange!50, draw=quotientorange!80!black] coordinates {(5,0.2)(6,0.1667)(7,0.1429)(8,0.125)};
\addplot[fill=targetgreen!45, draw=targetgreen!80!black] coordinates {(9,0.1111)(10,0.1)(11,0.0909)(12,0.0833)(13,0.0769)(14,0.0714)(15,0.0667)(16,0.0625)};
\draw[tangentred, thick, dashed] (axis cs:2.6,0.25) -- (axis cs:4.4,0.25);
\draw[tangentred, thick, dashed] (axis cs:4.6,0.125) -- (axis cs:8.4,0.125);
\draw[tangentred, thick, dashed] (axis cs:8.6,0.0625) -- (axis cs:16.4,0.0625);
\node[font=\small, curveblue] at (axis cs:3.5,0.48) {$2\cdot\tfrac14$};
\node[font=\small, quotientorange!80!black] at (axis cs:6.5,0.33) {$4\cdot\tfrac18$};
\node[font=\small, targetgreen!70!black] at (axis cs:12.5,0.24) {$8\cdot\tfrac1{16}$};
\end{axis}
\end{tikzpicture}
\caption{Oresme's grouping. Each colored block has twice as many terms as the last. Every bar in a block is at least as tall as its last bar (dashed red line), so each block adds at least $2\cdot\tfrac14 = 4\cdot\tfrac18 = 8\cdot\tfrac1{16} = \tfrac12$. Infinitely many blocks each adding $\tfrac12$ give an infinite total.}
\label{ch29:fig-oresme}
\end{figure}

\begin{connectionbox}
Nicole Oresme, a French scholar of the fourteenth century, gave this argument around 1350, centuries before calculus existed. It is one of the earliest rigorous results about an infinite process. The divergence is slow: to push the partial sum past $10$ you need over twelve thousand terms; to reach $100$ you need more terms than there are atoms in your body. This slowness is why the result surprises people: computed numerically, the partial sums look as if they are levelling off.
\end{connectionbox}

\begin{aestheticbox}
The harmonic series shows the main lesson of the chapter: \emph{what matters is how fast the terms shrink.} The terms $\tfrac1n$ go to zero, but not fast enough; the terms $\tfrac{1}{2^n}$ go to zero fast enough. Every test that follows asks this one question in a different form: \emph{do the terms shrink fast enough?} Geometric decay ($r^n$) is fast enough. Decay like $\tfrac1n$ is not. The borderline lies somewhere in between, and locating it is the job of the rest of the chapter.
\end{aestheticbox}

\section{The Integral Test: Series as Areas}
\label{ch29:integral-test}

The harmonic series raises a question: where, exactly, is the borderline? We know $\sum \tfrac1n$ diverges. What about $\sum \tfrac{1}{n^2}$, or $\sum \tfrac{1}{n^{1.5}}$, or $\sum \tfrac{1}{\sqrt n}$? To settle a whole family at once we connect sums to integrals: both measure area under a decreasing curve.

\begin{intuitionbox}
Picture the terms $a_n = f(n)$ of a positive, decreasing sequence as the heights of rectangles, each of width $1$. The sum $\sum f(n)$ is the total area of an infinite staircase of rectangles. The integral $\int_1^\infty f(x)\,\dd x$ is the area under the smooth curve $y=f(x)$. Because the curve is decreasing, the rectangles and the curve trap each other: the staircase is never far from the smooth area. So the sum is finite exactly when the integral is finite. We've traded a hard sum for an integral we may know how to do. Figure~\ref{ch29:fig-integraltest} shows the two ways the staircase and the curve compare.
\end{intuitionbox}

\begin{figure}[H]
\centering
\begin{tikzpicture}
% Panel (a): Left-endpoint (upper) rectangles: Sum dominates integral
\begin{axis}[
  name=upper,
  width=6.8cm, height=5.5cm,
  axis lines=left,
  xlabel={$x$}, ylabel={$y$},
  title style={font=\small},
  title={\textbf{(a)} Upper bound: $\int_1^N f \le \sum_{n=1}^{N-1} f(n)$},
  xtick={1,2,3,4,5}, ytick={0.5,1.0},
  xmin=0.5, xmax=5.5, ymin=0, ymax=1.2,
  domain=1:5, samples=80,
  clip=false,
]
% Upper rectangles
\addplot[ybar interval, fill=quotientorange, fill opacity=0.35, draw=quotientorange!80!black]
  coordinates {(1,1.0)(2,0.5)(3,0.333)(4,0.25)(5,0.2)};
\addplot[curveblue, very thick, domain=1:5] {1/x};
\node[quotientorange!80!black, font=\footnotesize, anchor=west] at (axis cs:2.2,0.85) {Rectangles $\ge$ Curve};
\end{axis}

% Panel (b): Right-endpoint (lower) rectangles: Integral bounds sum
\begin{axis}[
  name=lower,
  at={($(upper.east)+(1.0cm,0)$)},
  anchor=west,
  width=6.8cm, height=5.5cm,
  axis lines=left,
  xlabel={$x$}, ylabel={},
  title style={font=\small},
  title={\textbf{(b)} Lower bound: $\sum_{n=2}^N f(n) \le \int_1^N f$},
  xtick={1,2,3,4,5}, ytick=\empty,
  xmin=0.5, xmax=5.5, ymin=0, ymax=1.2,
  domain=1:5, samples=80,
  clip=false,
]
% Lower rectangles shifted by 1: on [n-1, n], height is f(n)
\addplot[ybar interval, fill=targetgreen, fill opacity=0.35, draw=targetgreen!80!black]
  coordinates {(1,0.5)(2,0.333)(3,0.25)(4,0.2)(5,0.167)};
\addplot[curveblue, very thick, domain=1:5] {1/x};
\node[targetgreen!80!black, font=\footnotesize, anchor=west] at (axis cs:2.2,0.85) {Rectangles $\le$ Curve};
\end{axis}

\end{tikzpicture}
\caption{The Integral Test sandwich. \textbf{(a)} Left-endpoint rectangles sit \emph{above} the decreasing curve $f(x)$, showing $\int_1^N f(x)\,\dd x \le \sum_{n=1}^{N-1} a_n$. If the integral diverges, it forces the series to diverge. \textbf{(b)} Right-endpoint rectangles sit \emph{below} the curve on $[n-1,n]$, showing $\sum_{n=2}^N a_n \le \int_1^N f(x)\,\dd x$. If the improper integral converges, it caps the series from above, guaranteeing convergence.}
\label{ch29:fig-integraltest}
\end{figure}

\begin{theorem}[The integral test]
\label{ch29:thm-integral}
Let $f$ be positive, continuous, and decreasing on $[1,\infty)$, and set $a_n = f(n)$. Then the series $\sum_{n=1}^{\infty} a_n$ and the improper integral $\int_1^{\infty} f(x)\,\dd x$ either both converge or both diverge.
\end{theorem}

\begin{proof}
Because $f$ decreases, on the interval $[n, n+1]$ we have $f(n+1) \le f(x) \le f(n)$. Integrating across that unit interval (width $1$) gives the rectangle-vs-curve sandwich:
\[
f(n+1) \;\le\; \int_{n}^{n+1} f(x)\,\dd x \;\le\; f(n).
\]
Add up the right-hand inequality for $n=1, 2, \dots, N-1$, and then do the same with the left-hand inequality:
\begin{align*}
\int_1^{N} f(x)\,\dd x &\le \sum_{n=1}^{N-1} f(n), &
\sum_{n=2}^{N} f(n) &\le \int_1^{N} f(x)\,\dd x.
\end{align*}
Together,
\[
\sum_{n=2}^{N} a_n \;\le\; \int_1^{N} f(x)\,\dd x \;\le\; \sum_{n=1}^{N-1} a_n .
\]
The partial sums increase, because every term is positive. If the integral converges (stays finite as $N\to\infty$), the left inequality bounds the partial sums above, and by the Monotone Convergence Theorem of Chapter~28 an increasing sequence that is bounded above converges; so the series converges. If the integral diverges, the right inequality forces the partial sums to infinity, so the series diverges.
\end{proof}

\begin{whybox}
The integral test works because the sum and the improper integral both measure area under the curve $y = f(x)$, one with rectangles and one exactly. When $f$ is decreasing, the gap between a rectangle and the curve on each unit interval $[n, n+1]$ is at most $f(n) - f(n+1)$. Adding these gaps gives a telescoping sum that is at most $f(1)$, so the total gap stays finite. This is why the test needs $f$ to be decreasing: that is what keeps the sum and the integral within a bounded distance of each other.
\end{whybox}

The payoff is a whole family of series settled at once. A \textbf{$p$-series} is a series $\sum 1/n^p$, where $p$ is a fixed real number; the harmonic series is the case $p=1$.

\begin{theorem}[The $p$-series]
\label{ch29:thm-pseries}
The series $\displaystyle\sum_{n=1}^{\infty} \frac{1}{n^p}$ converges if and only if $p > 1$.
\end{theorem}

\begin{proof}
Apply the integral test to $f(x)=x^{-p}$ (positive, decreasing for $x\ge 1$ when $p>0$; and for $p\le 0$ the terms don't even go to zero, so it diverges by the $n$th-term test). For $p \ne 1$,
\[
\int_1^{\infty} x^{-p}\,\dd x = \lim_{b\to\infty} \frac{x^{1-p}}{1-p}\Big|_1^{b} = \lim_{b\to\infty}\frac{b^{1-p}-1}{1-p}.
\]
When $p>1$, the exponent $1-p$ is negative, so $b^{1-p}\to 0$ and the integral equals $\tfrac{1}{p-1}$, finite: convergence. When $p<1$, $b^{1-p}\to\infty$: divergence. When $p=1$, $\int_1^\infty \tfrac1x\,\dd x = \lim_{b\to\infty}\ln b = \infty$: divergence (the harmonic series, again).
\end{proof}

\begin{focusbox}
\textbf{Memorize this.} $\sum \tfrac{1}{n^p}$ converges $\iff p > 1$. So $\sum\tfrac1{n^2}, \sum\tfrac1{n^3}, \sum\tfrac{1}{n^{1.01}}$ converge; $\sum\tfrac1n, \sum\tfrac{1}{\sqrt n}=\sum n^{-1/2}, \sum\tfrac{1}{n^{0.99}}$ diverge. The borderline you have been looking for since Section~\ref{ch29:harmonic} sits exactly at $p=1$, and the harmonic series is on the \emph{divergent} side of it: among the pure powers $1/n^p$, the terms $1/n$ shrink fastest of all those that still diverge. Figure~\ref{ch29:fig-pseries} shows the three cases side by side.
\end{focusbox}

\begin{figure}[H]
\centering
\begin{tikzpicture}
\begin{axis}[
    width=12cm, height=6cm,
    xlabel={$N$ (number of terms)}, ylabel={partial sum},
    xmin=0, xmax=20.5, ymin=0, ymax=3.7,
    xtick={2,5,10,15,20},
    legend style={at={(0.5,-0.28)}, anchor=north, font=\small}, legend cell align=left,
    grid=both, grid style={gray!18},
]
% p = 0.5 : divergent, partial sums of 1/sqrt(n)
\addplot[trapred, thick, mark=*, mark size=1.1pt]
    coordinates {(1,1)(2,1.7071)(3,2.2845)(4,2.7845)(5,3.2317)(6,3.6400)};
\addlegendentry{$p=1/2$ (diverges, climbs without bound)}
% p = 1 : harmonic, divergent but slow
\addplot[quotientorange, thick, mark=*, mark size=1.1pt]
    coordinates {(1,1)(2,1.5)(4,2.0833)(6,2.45)(8,2.7179)(10,2.9290)(14,3.2516)(18,3.4951)(20,3.5977)};
\addlegendentry{$p=1$ (harmonic, diverges slowly)}
% p = 2 : convergent, approaches pi^2/6 = 1.6449
\addplot[curveblue, thick, mark=*, mark size=1.1pt]
    coordinates {(1,1)(2,1.25)(3,1.3611)(4,1.4236)(6,1.4914)(8,1.5274)(10,1.5498)(14,1.5760)(20,1.5962)};
\addlegendentry{$p=2$ (converges to $\pi^2/6\approx1.645$)}
\addplot[curveblue, dashed, domain=0:20.5] {1.6449};
\end{axis}
\end{tikzpicture}
\caption{Partial sums of three $p$-series. For $p=2$ they flatten toward the dashed line at $\pi^2/6$; for $p=1$ and $p=\tfrac12$ they keep rising, the $p=\tfrac12$ case faster. Convergence depends only on whether $p > 1$.}
\label{ch29:fig-pseries}
\end{figure}

\begin{connectionbox}
For even whole numbers $p$, the sum of the $p$-series is a rational multiple of $\pi^p$; for example $\sum 1/n^4 = \pi^4/90$. Euler surprised Europe in 1735 by proving $\sum_{n=1}^\infty \tfrac{1}{n^2} = \tfrac{\pi^2}{6}$, a fact the integral test alone cannot deliver (it tells you \emph{that} the sum is finite, not \emph{what} it equals). Viewed as a function of $p$, the sum is called the Riemann zeta function, $\zeta(p) = \sum_{n=1}^\infty n^{-p}$ ($\zeta$ is the Greek letter zeta); so $\zeta(2) = \pi^2/6$ and $\zeta(3) = \sum 1/n^3$. Its behavior is closely tied to how the prime numbers are spread out among the whole numbers.
\end{connectionbox}

\section{Comparison Tests: Lean on a Series You Already Understand}
\label{ch29:comparison}

We now have two reliable yardsticks: geometric series and $p$-series. The comparison tests measure unfamiliar series against these known ones.

\begin{intuitionbox}
If your terms are smaller (term by term) than the terms of a series you already know converges, your sum is trapped underneath a finite total -- it has nowhere to go but to a finite value. And if your terms are bigger than those of a series you know diverges, your sum is pushed up past every bound. Smaller than finite is finite; bigger than infinite is infinite.
\end{intuitionbox}

\begin{theorem}[Direct comparison test]
\label{ch29:thm-comparison}
Suppose $0 \le a_n \le b_n$ for all $n$ (beyond some starting point).
\begin{itemize}
\item If $\sum b_n$ converges, then $\sum a_n$ converges.
\item If $\sum a_n$ diverges, then $\sum b_n$ diverges.
\end{itemize}
\end{theorem}

\begin{proof}
All terms are nonnegative, so both sequences of partial sums are increasing. Let $A_n=\sum_{k\le n} a_k$ and $B_n = \sum_{k\le n} b_k$; then $A_n \le B_n$. If $\sum b_n$ converges to $B$, then $A_n \le B_n \le B$, so $(A_n)$ is increasing and bounded above, hence convergent. The second bullet is the contrapositive of the first.
\end{proof}

\begin{example}
Does $\displaystyle\sum_{n=1}^\infty \frac{1}{n^2 + 3n + 1}$ converge? The denominator is bigger than $n^2$, so each term is \emph{smaller} than $\tfrac{1}{n^2}$:
\[
0 < \frac{1}{n^2+3n+1} < \frac{1}{n^2}.
\]
Since $\sum \tfrac1{n^2}$ converges ($p=2>1$), our smaller series converges too. We never needed its value, only that it is trapped under something finite. Figure~\ref{ch29:fig-comparison} shows the two sets of terms side by side.
\end{example}

\begin{figure}[htbp]
\centering
\begin{tikzpicture}
\begin{axis}[
    width=12cm, height=5.5cm,
    xlabel={$n$}, ylabel={term size},
    xmin=0.4, xmax=8.6, ymin=0, ymax=1.08,
    xtick={1,2,3,4,5,6,7,8},
    ybar, bar width=0.55, bar shift=0pt,
    axis lines=left,
    legend pos=north east, legend cell align=left,
    legend style={font=\small},
]
\addplot[fill=curveblue!25, draw=curveblue]
    coordinates {(1,1)(2,0.25)(3,0.1111)(4,0.0625)(5,0.04)(6,0.0278)(7,0.0204)(8,0.0156)};
\addlegendentry{$b_n = 1/n^2$ (sum is finite, $\pi^2/6$)}
\addplot[fill=targetgreen!70, draw=targetgreen!70!black, bar width=0.3]
    coordinates {(1,0.2)(2,0.0909)(3,0.0526)(4,0.0345)(5,0.0244)(6,0.0182)(7,0.0141)(8,0.0112)};
\addlegendentry{$a_n = 1/(n^2+3n+1)$}
\end{axis}
\end{tikzpicture}
\caption{Direct comparison. Each green bar $a_n$ fits inside the blue bar $b_n$ behind it. The blue bars have a finite total area, so the green bars, which use less area, have a finite total too.}
\label{ch29:fig-comparison}
\end{figure}

The direct comparison needs a term-by-term inequality, which can be fiddly to arrange (you may have to make the smaller series smaller, or the bigger one bigger, in just the right direction). Often it is easier to compare how the terms behave for large $n$, by looking at the ratio $a_n/b_n$. That is the limit comparison test.

\begin{theorem}[Limit comparison test]
\label{ch29:thm-limcomparison}
Let $a_n > 0$ and $b_n > 0$, and suppose
\[
\lim_{n\to\infty} \frac{a_n}{b_n} = L, \qquad 0 < L < \infty.
\]
Then $\sum a_n$ and $\sum b_n$ either both converge or both diverge.
\end{theorem}

\begin{whybox}
Why does a finite, nonzero limit force the two series to share a fate? Because $\tfrac{a_n}{b_n}\to L$ means that eventually $a_n$ is sandwiched between $\tfrac{L}{2}b_n$ and $2L\,b_n$ -- the two series are, in the long run, constant multiples of each other, up to a factor between $\tfrac L2$ and $2L$. Multiplying every term of a series by a fixed positive constant cannot change whether it converges (it only scales the total). Convergence is decided by the \emph{tail}, by how fast terms shrink, and a constant factor does not change how fast they shrink. So $\sum a_n$ and $\sum b_n$ have the same convergence behavior. You do not need an exact inequality; you only need to identify which known series your terms \emph{behave like}.
\end{whybox}

\begin{focusbox}
\textbf{How to pick $b_n$ in practice.} Strip your term down to its dominant behavior for large $n$: keep only the fastest-growing part of the numerator and of the denominator. What is left tells you the comparison series. Example: $a_n = \dfrac{3n^2 - n}{4n^4 + 5}$ behaves like $\dfrac{3n^2}{4n^4} = \dfrac{3}{4n^2}$, so compare with $b_n = \tfrac{1}{n^2}$.
\end{focusbox}

\begin{example}
Test $\displaystyle\sum_{n=1}^\infty \frac{3n^2 - n}{4n^4 + 5}$. The dominant behavior is $\tfrac{1}{n^2}$, so take $b_n = \tfrac1{n^2}$:
\[
\frac{a_n}{b_n} = \frac{3n^2 - n}{4n^4 + 5}\cdot n^2 = \frac{3n^4 - n^3}{4n^4 + 5} \longrightarrow \frac{3}{4}.
\]
The limit $\tfrac34$ is finite and nonzero, and $\sum\tfrac1{n^2}$ converges, so our series converges. This is quicker than finding a direct inequality.
\end{example}

\begin{example}[A divergence by limit comparison]
Test $\displaystyle\sum_{n=1}^\infty \frac{1}{\sqrt{n^2+1}}$. For large $n$, $\sqrt{n^2+1}\approx n$, so compare with $b_n=\tfrac1n$:
\[
\frac{a_n}{b_n} = \frac{n}{\sqrt{n^2+1}} = \frac{1}{\sqrt{1+1/n^2}} \longrightarrow 1.
\]
The limit is finite and nonzero. Since $\sum\tfrac1n$ diverges (harmonic), so does our series. The square root barely changed anything; the term still shrinks ``like $1/n$.''
\end{example}

\section{The Ratio and Root Tests: Asymptotically Geometric Series}
\label{ch29:ratio-root}

The geometric series converges because the ratio of consecutive terms is a constant $r$ with $\abs r < 1$. The \textbf{ratio test} applies this idea to series that are only \emph{approximately} geometric far out, with the approximation getting better as $n$ grows (\emph{asymptotically} geometric). This is a large class that includes most series built from factorials and exponentials.

\begin{intuitionbox}
Look at the size of the ratio $\abs{a_{n+1}/a_n}$ of each term to the previous one. For a geometric series $\sum a r^n$ this ratio is always $\abs{r}$. For a general series it varies, but if it \emph{settles} toward some limit $L$, then far out in the tail the series behaves like a geometric series with ratio $L$. So if $L<1$, the tail decays geometrically and the series converges; if $L>1$, the terms eventually \emph{grow}, and the series diverges.
\end{intuitionbox}

\begin{theorem}[Ratio test]
\label{ch29:thm-ratio}
For a series $\sum a_n$ with $a_n \neq 0$, suppose $\displaystyle\lim_{n\to\infty}\left|\frac{a_{n+1}}{a_n}\right| = L$.
\begin{itemize}
\item If $L < 1$, the series converges (in fact \emph{absolutely}, a stronger kind of convergence defined in Section~\ref{ch29:absolute}).
\item If $L > 1$ (or $L=\infty$), the series diverges.
\item If $L = 1$, the test is inconclusive -- it tells you nothing.
\end{itemize}
\end{theorem}

\begin{proof}[Why it works]
Suppose $L<1$. Pick a number $r$ strictly between $L$ and $1$ (for instance $r = \frac{L+1}{2}$). Applying the limit definition with tolerance $\eps = r - L > 0$ gives a cutoff $N$ with $\abs{a_{n+1}/a_n} < L + \eps = r$, that is $\abs{a_{n+1}} < r\,\abs{a_n}$, for all $n \ge N$. By induction $\abs{a_{N+k}} < r^k \abs{a_N}$, so the tail of the series is bounded, term by term, by the geometric series $\abs{a_N}\sum r^k$, which converges because $r<1$. By direct comparison the tail converges, hence the whole series converges absolutely. If instead $L>1$, then past some cutoff $\abs{a_{n+1}} > \abs{a_n}$, so the sizes of the terms keep growing, $a_n \not\to 0$, and the series diverges by the $n$th-term test. When $L=1$, the comparison geometric series sits exactly on its own boundary and tells us nothing.
\end{proof}

\begin{example}[Factorials suit the ratio test]
Test $\displaystyle\sum_{n=1}^\infty \frac{2^n}{n!}$. Form the ratio; most of the factorial and the power cancel:
\[
\left|\frac{a_{n+1}}{a_n}\right| = \frac{2^{n+1}}{(n+1)!}\cdot\frac{n!}{2^n} = \frac{2}{n+1} \longrightarrow 0.
\]
$L = 0 < 1$, so the series converges, and fast. (Its sum, as the next chapter shows, is $e^2 - 1$.) Whenever you see factorials, try the ratio test first; the cancellation is usually this clean. Figure~\ref{ch29:fig-ratio} shows the terms dropping under a geometric bound, which is why the test works.
\end{example}

\begin{example}[Why $L=1$ tells you nothing]
Apply the ratio test to $\sum \tfrac1n$ (diverges) and to $\sum \tfrac1{n^2}$ (converges):
\[
\sum\tfrac1n:\ \frac{1/(n+1)}{1/n}=\frac{n}{n+1}\to 1, \qquad
\sum\tfrac1{n^2}:\ \frac{n^2}{(n+1)^2}\to 1.
\]
The ratio test returns $L=1$ for both, yet one converges and one does not. So $L=1$ carries no information; you must switch to another test (here, the $p$-series rule). In particular the ratio test cannot decide any $p$-series, or any series whose terms shrink like a power $1/n^p$.
\end{example}

The \textbf{root test} is closely related to the ratio test and better suited when the $n$th term is an $n$th power. Instead of comparing neighbors, it extracts the ``geometric ratio'' directly by taking an $n$th root.

\begin{theorem}[Root test]
\label{ch29:thm-root}
For $\sum a_n$, suppose $\displaystyle\lim_{n\to\infty}\sqrt[n]{\abs{a_n}} = L$. If $L<1$ the series converges absolutely; if $L>1$ (or $\infty$) it diverges; if $L=1$ it is inconclusive.
\end{theorem}

\begin{whybox}
The root test works for the same reason as the ratio test, even more directly. If $\sqrt[n]{\abs{a_n}}\to L$ then for large $n$, $\abs{a_n}\approx L^n$: the term is roughly the $n$th term of a geometric series with ratio $L$. So $L<1$ means geometric decay and convergence; $L>1$ means geometric growth and divergence. The root test estimates the effective geometric ratio directly from $a_n$ itself, rather than from consecutive terms. Use it whenever $a_n$ involves an $n$th power, such as $a_n = \big(\tfrac{n}{2n+1}\big)^n$.
\end{whybox}

\begin{example}
Test $\displaystyle\sum_{n=1}^\infty \left(\frac{n}{2n+1}\right)^n$. Because the general term is an $n$th power, we apply the root test:
\[
\sqrt[n]{\abs{a_n}} = \frac{n}{2n+1} \longrightarrow \frac12 < 1,
\]
so the series converges. The ratio test would be messy here; the root test gives the answer in one line.
\end{example}

\begin{figure}[H]
\centering
\begin{tikzpicture}
\begin{axis}[
    width=12cm, height=6cm,
    xlabel={$n$}, ylabel={term size $|a_n|$},
    xmin=0.5, xmax=10.5, ymin=0, ymax=2.25,
    xtick={1,2,3,4,5,6,7,8,9,10},
    ytick={0.5,1.0,1.5,2.0},
    samples=60, domain=1:10.5,
    legend pos=north east, legend cell align=left,
    grid=both, grid style={gray!18},
]
% terms of sum 2^n/n!  : a_n = 2^n/n!
\addplot[only marks, mark=*, mark size=1.5pt, curveblue]
    coordinates {(1,2)(2,2)(3,1.3333)(4,0.6667)(5,0.2667)(6,0.0889)(7,0.0254)(8,0.00635)(9,0.001411)(10,0.000282)};
\addlegendentry{terms $2^n/n!$}
% geometric envelope past n=N: here once ratio < r=0.5, terms < a_5*(0.5)^{n-5}
\addplot[tangentred, dashed, very thick, domain=4:10.5] {0.6667*pow(0.5,x-4)};
\addlegendentry{geometric bound, ratio $\tfrac12$, from $n=4$}
\end{axis}
\end{tikzpicture}
\caption{The ratio test in action on $\sum 2^n/n!$. The ratio $2/(n+1)$ is below $\tfrac12$ once $n\ge 4$, so from there on every term lies on or under the dashed geometric sequence $a_4\left(\tfrac12\right)^{n-4}$, whose sum is finite. That is what forces convergence.}
\label{ch29:fig-ratio}
\end{figure}

\section{Alternating Series and the Sign That Saves You}
\label{ch29:alternating}

So far our terms have mostly been positive. When the signs alternate $+,-,+,-,\dots$, the partial cancellation can rescue a series that would diverge if all its terms were positive. The standard example is the \textbf{alternating harmonic series}:
\[
1 - \frac{1}{2} + \frac{1}{3} - \frac{1}{4} + \frac{1}{5} - \cdots = \sum_{n=1}^\infty \frac{(-1)^{n-1}}{n}.
\]
(The factor $(-1)^{n-1}$ is $+1$ for odd $n$ and $-1$ for even $n$; it supplies the alternating signs.) The all-positive version, the harmonic series, diverges. With alternating signs, this one \emph{converges}, to $\ln 2$ in fact.

\begin{intuitionbox}
Watch the partial sums of an alternating series whose terms shrink. You go up by $1$, down by $\tfrac12$, up by $\tfrac13$, down by $\tfrac14$. Each step overshoots the eventual limit, but by less than the last step (because the steps are shrinking). The partial sums zig-zag with ever-tightening swings, like a pendulum losing energy, homing in on a single value squeezed between every up-step and every down-step. Figure~\ref{ch29:fig-alternating} shows this for the alternating harmonic series.
\end{intuitionbox}

\begin{figure}[H]
\centering
\begin{tikzpicture}
\begin{axis}[
    width=12cm, height=6cm,
    xlabel={$n$}, ylabel={partial sum $S_n$},
    xmin=0.5, xmax=8.5, ymin=0.4, ymax=1.1,
    xtick={1,2,3,4,5,6,7,8},
    legend pos=north east, legend cell align=left,
    grid=both, grid style={gray!18},
]
\addplot[curveblue, thick, mark=*, mark size=1.5pt]
    coordinates {(1,1)(2,0.5)(3,0.8333)(4,0.5833)(5,0.7833)(6,0.6167)(7,0.7595)(8,0.6345)};
\addlegendentry{$S_n$ (alternating harmonic)}
\addplot[tangentred, dashed, thick, domain=0.5:8.5] {0.6931};
\addlegendentry{limit $\ln 2 \approx 0.693$}
\end{axis}
\end{tikzpicture}
\caption{Partial sums of $1 - \tfrac12 + \tfrac13 - \cdots$ oscillate with shrinking amplitude around $\ln 2$. The odd partial sums approach from above, the even ones from below, and they trap the limit between them.}
\label{ch29:fig-alternating}
\end{figure}

\begin{theorem}[Alternating series test, Leibniz]
\label{ch29:thm-alternating}
If the positive numbers $b_n$ satisfy (i) $b_{n+1} \le b_n$ (decreasing) and (ii) $b_n \to 0$, then the alternating series $\sum_{n=1}^\infty (-1)^{n-1} b_n$ converges. In addition, the error from stopping at the $n$th term is bounded by the next term: $\abs{S - S_n} \le b_{n+1}$.
\end{theorem}

\begin{proof}
First look at the even partial sums. Grouping in pairs,
\[
S_{2n} = (b_1 - b_2) + (b_3 - b_4) + \cdots + (b_{2n-1}-b_{2n});
\]
each pair is $\ge 0$ since $b_k$ decreases, so $(S_{2n})$ is increasing. Now group differently, starting one slot later:
\[
S_{2n} = b_1 - (b_2 - b_3) - (b_4 - b_5) - \cdots - (b_{2n-2} - b_{2n-1}) - b_{2n};
\]
every bracketed difference is $\ge 0$ and $b_{2n} > 0$, so $S_{2n} \le b_1$. Being increasing and bounded above, $(S_{2n})$ converges to a limit $S$ (Monotone Convergence Theorem). Meanwhile $S_{2n+1} = S_{2n} + b_{2n+1} \to S + 0 = S$, and $(S_{2n+1})$ decreases, since $S_{2n+1} - S_{2n-1} = -(b_{2n} - b_{2n+1}) \le 0$. So all partial sums converge to $S$, and $S_{2k} \le S \le S_{2k+1}$ for every $k$: the true sum always sits between two consecutive partial sums $S_n$ and $S_{n+1}$, which differ by exactly $b_{n+1}$. Hence $\abs{S - S_n} \le b_{n+1}$.
\end{proof}

\begin{aestheticbox}
The alternating series error bound is easy to use. To approximate the sum to within $0.001$, add terms until the next term is below $0.001$; the error from stopping there (the \emph{truncation error}) is smaller than that first omitted term. Because the error estimate needs no extra calculation, alternating series are convenient in numerical work.
\end{aestheticbox}

\section{Absolute vs.\ Conditional Convergence: A Subtle Hierarchy}
\label{ch29:absolute}

The alternating harmonic series converges, but only thanks to its signs: strip them away and it diverges. Contrast that with $\sum \tfrac{(-1)^{n-1}}{n^2}$, which converges even with all signs made positive (it's the $p$-series with $p=2$). These two convergences are not the same \emph{kind}, and the difference matters, as this section shows.

\begin{definition}[Absolute and conditional convergence]
\label{ch29:def-absolute}
A series $\sum a_n$ \textbf{converges absolutely} if $\sum \abs{a_n}$ converges. It \textbf{converges conditionally} if $\sum a_n$ converges but $\sum \abs{a_n}$ diverges -- that is, the convergence depends on cancellation among the signs.
\end{definition}

\begin{theorem}[Absolute convergence implies convergence]
\label{ch29:thm-abs-implies}
If $\sum \abs{a_n}$ converges, then $\sum a_n$ converges.
\end{theorem}

\begin{proof}
Each $a_n + \abs{a_n}$ is either $0$ (when $a_n<0$) or $2\abs{a_n}$ (when $a_n \ge 0$), so $0 \le a_n + \abs{a_n} \le 2\abs{a_n}$. Since $\sum 2\abs{a_n}$ converges, by direct comparison $\sum (a_n + \abs{a_n})$ converges. Then $\sum a_n = \sum (a_n + \abs{a_n}) - \sum \abs{a_n}$ is a difference of two convergent series, hence converges.
\end{proof}

\begin{focusbox}
This theorem is what makes the ratio and root tests valid for series with mixed signs. When they report ``$L<1$,'' they are proving that $\sum\abs{a_n}$ converges (absolute convergence), and by the theorem that guarantees ordinary convergence. So those tests handle signs automatically: you can apply them to a sign-changing series, because they only look at the sizes $\abs{a_n}$.
\end{focusbox}

\begin{happeningbox}
What is the structural difference between the two kinds? A \emph{rearrangement} of a series adds exactly the same terms, each used once, but in a different order. When $\sum\abs{a_n}$ converges, the positive terms alone have some finite total $P$ and the negative terms alone have some finite total $-Q$, so every rearrangement still adds up to $P - Q$: the order does not matter, exactly as with a finite sum. When a series converges only \emph{conditionally}, its positive terms alone sum to $+\infty$ and its negative terms alone sum to $-\infty$ (if either total were finite, convergence of the series would force the other to be finite too, making the convergence absolute). That gives you two unlimited supplies of ever-shrinking steps. Pick \emph{any} target $T$, say $T = 100$. Add positive terms in order until the running total first passes $100$ (it must, since that supply is unlimited); then add negative terms until it first dips below $100$; then positive again, back and forth forever. Because the individual terms shrink to $0$, the overshoots on either side of $100$ shrink to $0$ too, and the rearranged series converges to $100$. Figure~\ref{ch29:fig-rearrange} carries this out for the alternating harmonic series with target $T = 1.5$. This is the Riemann Rearrangement Theorem: absolute convergence makes a sum independent of order, while a conditionally convergent sum is a balance between two infinite supplies, and changing the order changes the balance.
\end{happeningbox}

\begin{figure}[htbp]
\centering
\begin{tikzpicture}
\begin{axis}[
    width=13cm, height=6cm,
    xlabel={number of terms used}, ylabel={running total},
    xmin=0, xmax=41, ymin=0.5, ymax=1.7,
    xtick={1,5,10,15,20,25,30,35,40},
    ytick={0.6931,1.0,1.5},
    yticklabels={$\ln 2$,$1$,$1.5$},
    ticklabel style={font=\small},
    legend style={at={(0.98,0.22)}, anchor=south east}, legend cell align=left,
    legend style={font=\small},
    grid=both, grid style={gray!18},
]
\addplot[curveblue, thick, mark=*, mark size=1.1pt]
    coordinates {(1,1.0000)(2,1.3333)(3,1.5333)(4,1.0333)(5,1.1762)(6,1.2873)(7,1.3782)(8,1.4551)(9,1.5218)(10,1.2718)(11,1.3306)(12,1.3833)(13,1.4309)(14,1.4744)(15,1.5144)(16,1.3477)(17,1.3847)(18,1.4192)(19,1.4515)(20,1.4818)(21,1.5103)(22,1.3853)(23,1.4124)(24,1.4380)(25,1.4624)(26,1.4857)(27,1.5079)(28,1.4079)(29,1.4292)(30,1.4496)(31,1.4692)(32,1.4880)(33,1.5062)(34,1.4229)(35,1.4404)(36,1.4574)(37,1.4738)(38,1.4896)(39,1.5050)(40,1.4336)};
\addlegendentry{rearranged running total}
\addplot[tangentred, dashed, thick, domain=0:41] {1.5};
\addlegendentry{target $T = 1.5$}
\addplot[gray, dotted, thick, domain=0:41] {0.6931};
\addlegendentry{original sum $\ln 2$}
\end{axis}
\end{tikzpicture}
\caption{Rearranging $1 - \tfrac12 + \tfrac13 - \tfrac14 + \cdots$ to aim at $1.5$. Positive terms $1, \tfrac13, \tfrac15, \dots$ are added until the total passes $1.5$, then one negative term ($-\tfrac12$, then $-\tfrac14$, \dots) pulls it back below, and so on. The swings shrink and the running total settles on $1.5$, not on $\ln 2$, even though every term of the original series is used exactly once.}
\label{ch29:fig-rearrange}
\end{figure}

\begin{connectionbox}
In 1837, \emph{Peter Gustav Lejeune Dirichlet} showed that infinite sums do not always obey the commutative law, the rule that the order of addition does not matter. While $a+b=b+a$ holds for finite sums, Dirichlet noticed that the alternating harmonic series $\sum_{n=1}^\infty \frac{(-1)^{n-1}}{n} = \ln 2$ could be rearranged (one positive term, then two negative terms, repeated) to sum to half its original value:
\begin{align*}
  1 - \frac12 + \frac13 - \frac14 + \cdots &= \ln 2, \qquad \text{but}\\
  \left(1 - \frac12 - \frac14\right) + \left(\frac13 - \frac16 - \frac18\right) + \cdots &= \frac12 \ln 2.
\end{align*}
Mathematicians had long assumed that an infinite sum would obey the commutative law; Dirichlet proved that it fails for conditionally convergent series.

In 1854, \emph{Bernhard Riemann} generalized this in the \textbf{Riemann Rearrangement Theorem}: if a series converges conditionally, its terms can be reordered to converge to any chosen real number, or to diverge to $\pm\infty$. Conditional convergence depends on the specific order of summation; only \emph{absolute convergence} guarantees a sum that does not depend on the order.
\end{connectionbox}

\section{Choosing a Test: A Systematic Strategy}
\label{ch29:strategy}

With this many tests available, the skill is knowing which to use. The \emph{form} of the term usually tells you. The list below gives a good order to try them in, and Figure~\ref{ch29:fig-strategy} draws it as a flowchart.

\begin{focusbox}
\textbf{A practical order of attack.}
\begin{enumerate}
\item \textbf{Glance at the terms.} Do they fail to go to $0$? Then it diverges by the $n$th-term test (Sec.~\ref{ch29:nth-term}). This takes one look, so always do it first.
\item \textbf{Recognize the type.} Is it geometric ($a r^n$) or telescoping (a difference $b_k - b_{k+1}$)? Then sum it \emph{exactly} (Secs.~\ref{ch29:geometric},~\ref{ch29:telescoping}).
\item \textbf{Spot factorials or $n$th powers.} Factorials or exponentials $\Rightarrow$ ratio test. A whole term raised to the $n$th power $\Rightarrow$ root test (Sec.~\ref{ch29:ratio-root}).
\item \textbf{See a fraction built from powers and roots of $n$:} find the power of $n$ it behaves like for large $n$ and use limit comparison against a $p$-series (Secs.~\ref{ch29:comparison},~\ref{ch29:integral-test}).
\item \textbf{Notice alternating signs} $(-1)^n$: check absolute convergence first (it's stronger and often easier); if that fails, try the alternating series test (Secs.~\ref{ch29:alternating},~\ref{ch29:absolute}).
\item \textbf{If $f(n)$ has an antiderivative you can find} ($\tfrac1{n\ln n}$, $n e^{-n^2}$): the integral test may be the most direct route (Sec.~\ref{ch29:integral-test}).
\end{enumerate}
\end{focusbox}

\begin{figure}[htbp]
\centering
\begin{tikzpicture}[
  font=\small,
  q/.style={draw=curveblue, fill=curveblue!8, rounded corners, align=center, minimum width=5.2cm, minimum height=0.9cm},
  a/.style={draw=targetgreen!70!black, fill=targetgreen!10, rounded corners, align=center, minimum width=4.8cm, minimum height=0.9cm},
  d/.style={draw=trapred, fill=trapred!8, rounded corners, align=center, minimum width=4.8cm, minimum height=0.9cm},
  arr/.style={->, thick, >=stealth},
  node distance=0.55cm and 0.9cm,
]
\node[q] (q1) {Do the terms $a_n$ go to $0$?};
\node[d, right=of q1] (r1) {Diverges ($n$th-term test)};
\node[q, below=of q1] (q2) {Geometric or telescoping?};
\node[a, right=of q2] (r2) {Sum it exactly};
\node[q, below=of q2] (q3) {Factorials, exponentials,\\ or an $n$th power?};
\node[a, right=of q3] (r3) {Ratio test or root test};
\node[q, below=of q3] (q4) {Fraction made of powers\\ and roots of $n$?};
\node[a, right=of q4] (r4) {Limit comparison\\ with a $p$-series};
\node[q, below=of q4] (q5) {Alternating signs?};
\node[a, right=of q5] (r5) {Test $\sum\abs{a_n}$ first; if that\\ diverges, alternating series test};
\node[a, below=of q5, minimum width=5.2cm] (r6) {Try the integral test,\\ or compare directly};
\draw[arr] (q1) -- node[above, font=\footnotesize] {no} (r1);
\draw[arr] (q1) -- node[left, font=\footnotesize] {yes} (q2);
\draw[arr] (q2) -- node[above, font=\footnotesize] {yes} (r2);
\draw[arr] (q2) -- node[left, font=\footnotesize] {no} (q3);
\draw[arr] (q3) -- node[above, font=\footnotesize] {yes} (r3);
\draw[arr] (q3) -- node[left, font=\footnotesize] {no} (q4);
\draw[arr] (q4) -- node[above, font=\footnotesize] {yes} (r4);
\draw[arr] (q4) -- node[left, font=\footnotesize] {no} (q5);
\draw[arr] (q5) -- node[above, font=\footnotesize] {yes} (r5);
\draw[arr] (q5) -- node[left, font=\footnotesize] {no} (r6);
\end{tikzpicture}
\caption{The order of attack as a flowchart. Start at the top and stop at the first box on the right that applies. The $n$th-term question comes first because it is the quickest.}
\label{ch29:fig-strategy}
\end{figure}

\begin{trapbox}
The classic mistakes:
\begin{itemize}
\item Concluding \emph{convergence} from $a_n\to 0$. The $n$th-term test only proves divergence; it is silent on convergence. (Harmonic series.)
\item Drawing a conclusion when the ratio or root test returns $L=1$. That result is \emph{inconclusive}: switch tests, do not guess.
\item Using the integral test when $f$ is not eventually decreasing. Without that hypothesis the rectangle-curve sandwich falls apart.
\item Using limit comparison with $L=0$ or $L=\infty$ carelessly. The ``both converge or both diverge'' conclusion needs $0<L<\infty$; the boundary cases give only one-directional information.
\item Reordering a conditionally convergent series as if order didn't matter. It does (Riemann).
\end{itemize}
\end{trapbox}

\begin{labbox}
Classify each by hand, naming the test and the reasoning before you compute. (a) $\sum \dfrac{n}{n^2+1}$. \;(b) $\sum \dfrac{n!}{n^n}$. \;(c) $\sum \dfrac{(-1)^n}{\ln n}$ (from $n=2$). \;(d) $\sum \dfrac{\cos n}{n^2}$.

\emph{Hints to check yourself:} (a) limit-compare with $\tfrac1n$, diverges. (b) ratio test, $\big(\tfrac{n}{n+1}\big)^n \to \tfrac1e<1$, converges. (c) terms decrease to $0$, alternating series test, converges -- but $\sum\tfrac{1}{\ln n}$ diverges (since $\ln n < n$ gives $\tfrac{1}{\ln n} > \tfrac1n$), so only \emph{conditionally}. (d) $\big|\tfrac{\cos n}{n^2}\big|\le \tfrac{1}{n^2}$, comparison gives \emph{absolute} convergence.
\end{labbox}

\begin{applicationbox}
Series arise throughout applied mathematics and scientific computing. In physics, the electrical potential energy of a crystal is an infinite (often conditionally convergent) sum over its ions (electrically charged atoms), so the order of summation must be specified. In probability, the expected value of a random variable taking values $x_k$ with probabilities $p_k$ is the series $\sum x_k p_k$, and the expected value is only considered well-defined when this series converges \emph{absolutely}. In reinforcement learning (training a computer program by rewarding good actions), a reward $r_t$ received $t$ steps in the future is counted as $\gamma^t r_t$, where the \emph{discount factor} $\gamma$ satisfies $0 \le \gamma < 1$; if the rewards are bounded, the total $\sum_{t=0}^\infty \gamma^t r_t$ is finite by comparison with a geometric series. And special functions (standard named functions beyond the elementary ones, such as the error function and $\zeta$) are evaluated numerically as series, whose convergence rate determines how many terms are needed for a given accuracy. The tests in this chapter guarantee that these computations make sense.
\end{applicationbox}

\section{Scope and Limits of Convergence Tests}

\begin{connectionbox}
Before Cauchy's 1821 definition of convergence through partial sums, infinite series often produced contradictory values (as with Grandi and Euler and $1 - 1 + 1 - \cdots$). The tests in this chapter serve two main purposes:
\begin{itemize}
  \item \emph{Certifying that a quantity exists:} proving that constants such as the
  Euler--Mascheroni constant $\gamma$ or $\zeta(3)$, and special functions such as the
  zeta function, have well-defined finite values, even when no formula for them is known;
  \item \emph{Computing with guaranteed accuracy:} error bounds such as the alternating
  series bound $|S - S_n| \le b_{n+1}$ and the geometric tail formula tell you how many terms
  are enough.
\end{itemize}

The tests also have two limits:
\begin{enumerate}
  \item \emph{They rarely give the value.} Most tests (integral test, comparison, ratio, root)
  only decide whether a series converges, not what its sum is. Proving that
  $\sum_{n=1}^\infty \frac{1}{n^2}$ converges takes half a minute using the integral test. Finding its
  exact sum, $\frac{\pi^2}{6}$ (the ``Basel problem'', solved by Euler in 1735), required
  different methods.
  \item \emph{Conditional convergence depends on order.} As Section~\ref{ch29:absolute}
  showed, a conditionally convergent sum can be changed by reordering its terms.
\end{enumerate}

In Chapter~30 (\emph{Power Series and Taylor Series}), we let the terms of a series
depend on a variable $x$:
\[
  f(x) = \sum_{n=0}^{\infty} c_n (x - a)^n ,
\]
where the $c_n$ are fixed numbers (the coefficients) and $a$ is a fixed point.
Treating such a series as an infinitely long polynomial lets us differentiate and integrate
term by term, write functions like $e^x$, $\sin x$, and $\arctan x$ as explicit power series,
find exact sums of numerical series, and solve differential equations.
\end{connectionbox}

\section{Exercises}
\label{ch29:exercises}

\begin{exercise}[Sum it exactly]
Evaluate each convergent series in closed form.
(a) $\displaystyle\sum_{n=0}^\infty \frac{3}{4^n}$. \quad
(b) $\displaystyle\sum_{n=1}^\infty \left(\frac{2}{3}\right)^{n}$. \quad
(c) $\displaystyle\sum_{n=1}^\infty \frac{1}{(2n-1)(2n+1)}$ (telescope via $\tfrac12\big(\tfrac{1}{2n-1}-\tfrac{1}{2n+1}\big)$). \quad
(d) Express $0.1\overline{6}$ (that is, $0.16666\ldots$) as a fraction using a geometric series.
\end{exercise}

\begin{exercise}[The $n$th-term test, used correctly]
Determine which of these diverge by the $n$th-term test, and explain why the test is \emph{silent} on the others.
(a) $\sum \dfrac{n^2}{n^2+5}$. \;
(b) $\sum \dfrac{1}{n^2+5}$. \;
(c) $\sum \cos\!\big(\tfrac1n\big)$. \;
(d) $\sum (-1)^n \dfrac{n}{n+1}$.
\end{exercise}

\begin{exercise}[$p$-series and comparison]
Decide convergence, naming your comparison series each time.
(a) $\sum \dfrac{1}{n^{3/2}}$. \;
(b) $\sum \dfrac{1}{\sqrt[3]{n}}$. \;
(c) $\sum \dfrac{n+4}{n^3 - 2n + 7}$. \;
(d) $\sum \dfrac{1}{\sqrt{n^3+1}}$.
\end{exercise}

\begin{exercise}[Ratio and root]
Apply the ratio or root test, whichever is cleaner, and say what $L$ is.
(a) $\sum \dfrac{5^n}{n!}$. \;
(b) $\sum \dfrac{n!}{10^n}$. \;
(c) $\sum \dfrac{n^2}{2^n}$. \;
(d) $\sum \left(\dfrac{3n+1}{4n+5}\right)^{n}$.
\end{exercise}

\begin{exercise}[Integral test]
Use the integral test to settle each, and verify the hypotheses (positive, decreasing) hold.
(a) $\displaystyle\sum_{n=2}^\infty \frac{1}{n\ln n}$. \;
(b) $\displaystyle\sum_{n=2}^\infty \frac{1}{n(\ln n)^2}$. \;
(c) $\displaystyle\sum_{n=1}^\infty n\,e^{-n^2}$.
Comment on how (a) and (b) differ despite looking nearly identical -- where is the borderline now?
\end{exercise}

\begin{exercise}[Alternating, absolute, conditional]
For each, decide whether it converges absolutely, converges conditionally, or diverges. Justify all three possibilities you rule out or confirm.
(a) $\displaystyle\sum_{n=1}^\infty \frac{(-1)^{n-1}}{\sqrt n}$. \;
(b) $\displaystyle\sum_{n=1}^\infty \frac{(-1)^{n-1}}{n^2}$. \;
(c) $\displaystyle\sum_{n=1}^\infty \frac{(-1)^{n-1}\,n}{n+1}$. \;
(d) $\displaystyle\sum_{n=2}^\infty \frac{(-1)^n}{n\ln n}$.
\end{exercise}

\begin{exercise}[Error estimate]
For the alternating series $\displaystyle\sum_{n=1}^\infty \frac{(-1)^{n-1}}{n^3}$, how many terms must you add to guarantee the partial sum is within $10^{-3}$ of the true value? Use the alternating-series error bound, and explain why that bound applies here.
\end{exercise}

\begin{exercise}[Conceptual]
Answer in a sentence or two each, with reasons.
(a) True or false: if $a_n \to 0$ then $\sum a_n$ converges. Give the canonical counterexample.
(b) True or false: if $\sum a_n$ converges then $a_n \to 0$. Which theorem is this?
(c) Can a series of positive terms converge conditionally? Why or why not?
(d) If $\sum a_n$ and $\sum b_n$ both diverge, must $\sum(a_n+b_n)$ diverge? Find an example or prove it.
\end{exercise}

\begin{exercise}[Prove it]
Prove that if $\sum a_n$ converges and $\sum b_n$ diverges, then $\sum (a_n + b_n)$ diverges. (Hint: argue by contradiction using the fact that the difference of two convergent series converges.)
\end{exercise}

\begin{exercise}[Prove it: the $n$th-term test]
Prove the $n$th-term test from scratch: if $\sum_{k=1}^\infty a_k = S$ (converges), then $a_n \to 0$. Your proof should use only the definition of convergence via partial sums (Definition~\ref{ch29:def-convergence}) and the fact that if $S_n \to S$ then $S_{n-1}\to S$. Then explain in one paragraph, in your own words, why the converse fails, citing a specific series.
\end{exercise}

\begin{exercise}[Drug Accumulation as a Geometric Series]
A patient takes a $100$~mg dose of a drug each day. Each day, the body eliminates $25\%$ of whatever drug is present, so $75\%$ carries over to the next day. Model the long-run amount of drug in the body just after a dose as an infinite series, identify it as geometric, and find the steady-state amount. (This shares the discounted-sum structure discussed in the application box.)
\end{exercise}

\begin{exercise}[Build the borderline]
We saw $\sum \tfrac1n$ diverges and $\sum\tfrac1{n^2}$ converges. Investigate the finer scale between them:
(a) Argue that $\sum \tfrac{1}{n^{1.0001}}$ converges and $\sum\tfrac{1}{n^{0.9999}}$ diverges -- the borderline at $p=1$ is sharp.
(b) Using the integral test, show $\sum \tfrac{1}{n\ln n}$ still diverges, so $\tfrac1n$ is not even the ``last'' divergent term -- there is a whole ladder of ever-slower divergences below it. Reflect briefly on what this says about the idea of a single ``boundary'' between convergence and divergence.
\end{exercise}
